% GENERATED FILE. Edit canonical lessons, then run npm run generate:books.
% canonical-source-hash: fnv1a64:5167a39aa4da6f01
% canonical-lessons: ML-C64-then, ML-C64-soon, ML-C64-sometimes, ML-C64-only, ML-C64-but

\chapter{Five Words That Join}
\label{ch:ml-five-words-that-join}
\hlchaptermodality{64}

\begin{chapteropening}
\textbf{By the end of this chapter:} \emph{I can say then, soon, sometimes, only and but in Malayalam.}

Complete ML-C64-but and say all five words in order.
\end{chapteropening}

\section[\texorpdfstring{\ml{പിന്നെ} (pinne)}{പിന്നെ (pinne)}]{\ml{പിന്നെ} (pinne) — then, after that}

\label{lesson:ML-C64-then}



\begin{quote}
Before the new run of words: what did \emph{cuma} mean, and what did \emph{pani} mean?
\end{quote}

\subsection*{You'll want to know: \ml{പിന്നെ}}
\textbf{\ml{പിന്നെ}} (\emph{pinne}) — \textquotedblleft{}then, after that\textquotedblright{}.

Native Dravidian, built on \emph{pin}, the piece meaning behind or at the back. What sits behind you in space becomes what comes after you in time, and the same short word does both jobs.

Kerala also leans on it as a filler, the \emph{so...} a speaker puts in front of a sentence to buy a second before starting it.

After \emph{śari}, it turns an agreement into a continuation: right, and then.

\begin{cousinweb}
\textbf{\ml{പിന്നെ}} (\emph{pinne}) is native Dravidian, built on \emph{pin}, \textquotedblleft{}behind, after\textquotedblright{}; DEDR 4205 lists the root with Tamil \emph{piṉ}.
\end{cousinweb}

\begin{grammarlens}[title={What you've built}]
The first of five small words that join one thing to the next.
\end{grammarlens}

\subsection*{Your turn}
\begin{itemize}
\raggedright
  \item \emph{Say it:} \emph{pinne}
  \item \emph{Say it:} it once more, slowly
  \item \emph{Say it:} \emph{śari}, then \emph{pinne}, as though agreeing and going on
  \item \emph{Recall:} say \emph{kayaṟ}
\end{itemize}

\begin{tcolorbox}[breakable,colback=teal!4,colframe=teal!35!black,title={Before you move on}]
What does \ml{പിന്നെ} mean? (\textquotedblleft{}then, after that\textquotedblright{}.) Say it once more.
\end{tcolorbox}
\section[\texorpdfstring{\ml{ഉടനെ} (uṭane)}{ഉടനെ (uṭane)}]{\ml{ഉടനെ} (uṭane) — soon, right away}

\label{lesson:ML-C64-soon}



\begin{quote}
Before the new one: what did \emph{pinne} mean?
\end{quote}

\subsection*{You'll want to know: \ml{ഉടനെ}}
\textbf{\ml{ഉടനെ}} (\emph{uṭane}) — \textquotedblleft{}soon, right away\textquotedblright{}.

Native Dravidian, built on \emph{uṭan}, meaning along with — what happens beside a thing rather than after it.

So the two words you now have divide the near future between them. \emph{Pinne} leaves a gap; \ml{ഉടനെ} closes it.

In Kerala speech it also softens a refusal: not now, but not far off either.

\begin{cousinweb}
\textbf{\ml{ഉടനെ}} (\emph{uṭane}) is native Dravidian, built on \emph{uṭan}, \textquotedblleft{}together with, at once\textquotedblright{}; DEDR 945 lists the root with Tamil \emph{uṭaṉ}.
\end{cousinweb}

\begin{grammarlens}[title={What you've built}]
Two.
\end{grammarlens}

\subsection*{Your turn}
\begin{itemize}
\raggedright
  \item \emph{Say it:} \emph{uṭane}
  \item \emph{Say it:} it once more, slowly
  \item \emph{Say it:} \emph{pinne}, then \emph{uṭane}, and say which of the two leaves a gap
  \item \emph{Read it:} \textbf{\ml{നൂല്}}
\end{itemize}

\begin{tcolorbox}[breakable,colback=teal!4,colframe=teal!35!black,title={Before you move on}]
What does \ml{ഉടനെ} mean? (\textquotedblleft{}soon, right away\textquotedblright{}.) Say it once more.
\end{tcolorbox}
\section[\texorpdfstring{\ml{ചിലപ്പോൾ} (cilappōḷ)}{ചിലപ്പോൾ (cilappōḷ)}]{\ml{ചിലപ്പോൾ} (cilappōḷ) — sometimes}

\label{lesson:ML-C64-sometimes}



\begin{quote}
Before the new one: what did \emph{uṭane} mean?
\end{quote}

\subsection*{You'll want to know: \ml{ചിലപ്പോൾ}}
\textbf{\ml{ചിലപ്പോൾ}} (\emph{cilappōḷ}) — \textquotedblleft{}sometimes\textquotedblright{}.

Two pieces, and you can take them apart by ear: \emph{cila}, some or a few, and \emph{pōḷ}, the when-piece.

That tail is the same one closing \emph{ippōḷ}, the now-word. One says at this time, the other says at some times, and the shared ending is the reason the two rhyme.

It is the word a Malayali reaches for when hedging: not a denial, only a fence around a claim.

\begin{cousinweb}
\textbf{\ml{ചിലപ്പോൾ}} (\emph{cilappōḷ}) is native Dravidian, \textquotedblleft{}some times\textquotedblright{}: \emph{cila}, \textquotedblleft{}some\textquotedblright{} (DEDR 1571, with Tamil \emph{cila}), joined to a word for \textquotedblleft{}time\textquotedblright{} (DEDR 4559, with Tamil \emph{pōtu}), the same word that ends \emph{ippōḷ}.
\end{cousinweb}

\begin{grammarlens}[title={What you've built}]
Three.
\end{grammarlens}

\subsection*{Your turn}
\begin{itemize}
\raggedright
  \item \emph{Say it:} \emph{cilappōḷ}
  \item \emph{Say it:} it once more, slowly
  \item \emph{Say it:} \emph{ippōḷ}, then \emph{cilappōḷ}, and say which of the two is certain
  \item \emph{Recall:} say \emph{sūci}
\end{itemize}

\begin{tcolorbox}[breakable,colback=teal!4,colframe=teal!35!black,title={Before you move on}]
What does \ml{ചിലപ്പോൾ} mean? (\textquotedblleft{}sometimes\textquotedblright{}.) Say it once more.
\end{tcolorbox}
\section[\texorpdfstring{\ml{മാത്രം} (mātraṁ)}{മാത്രം (mātraṁ)}]{\ml{മാത്രം} (mātraṁ) — only, that much and no more}

\label{lesson:ML-C64-only}



\begin{quote}
Before the new one: what did \emph{cilappōḷ} mean?
\end{quote}

\subsection*{You'll want to know: \ml{മാത്രം}}
\textbf{\ml{മാത്രം}} (\emph{mātraṁ}) — \textquotedblleft{}only, that much and no more\textquotedblright{}.

Sanskrit \emph{mātra}, a measure, from the root \emph{mā}, to measure. The Greek \emph{metron} behind \emph{metre} grows from the same seed, so this word and a poet's line-length are cousins.

Set after a word it fences that word off: this much, and nothing beyond it.

Beside \emph{kuṟaccŭ}, which says a small amount, \ml{മാത്രം} says an exact amount and shuts the door.

\begin{cousinweb}
\textbf{\ml{മാത്രം}} (\emph{mātraṁ}) $\leftarrow$ Sanskrit \emph{mātrā}, \textquotedblleft{}measure\textquotedblright{}, taken over whole as a tatsama (CDIAL 10023). At the end of a Sanskrit compound \emph{-mātra} already meant \textquotedblleft{}only so much\textquotedblright{}.
\end{cousinweb}

\begin{grammarlens}[title={What you've built}]
Four.
\end{grammarlens}

\subsection*{Your turn}
\begin{itemize}
\raggedright
  \item \emph{Say it:} \emph{mātraṁ}
  \item \emph{Say it:} it once more, slowly
  \item \emph{Say it:} \emph{kuṟaccŭ}, then \emph{mātraṁ}, and say which one sets a limit
  \item \emph{Read it:} \textbf{\ml{ചൂല്}}
\end{itemize}

\begin{tcolorbox}[breakable,colback=teal!4,colframe=teal!35!black,title={Before you move on}]
What does \ml{മാത്രം} mean? (\textquotedblleft{}only, that much and no more\textquotedblright{}.) Say it once more.
\end{tcolorbox}
\section[\texorpdfstring{\ml{എന്നാൽ} (ennāl)}{എന്നാൽ (ennāl)}]{\ml{എന്നാൽ} (ennāl) — but, on the other hand}

\label{lesson:ML-C64-but}



\begin{quote}
Before the new one: what did \emph{mātraṁ} mean?
\end{quote}

\subsection*{You'll want to know: \ml{എന്നാൽ}}
\textbf{\ml{എന്നാൽ}} (\emph{ennāl}) — \textquotedblleft{}but, on the other hand\textquotedblright{}.

Built out of the saying-piece \emph{enn-}, the one behind reported speech, and a tail meaning by or through: roughly, by that having been said.

English arrives at almost the same phrase from the other end with \emph{that said}, which is a small comfort — the shape is already familiar even though the pieces are not.

So the word grants the previous sentence before it turns on it, which is a gentler move than a flat contradiction.

\begin{cousinweb}
\textbf{\ml{എന്നാൽ}} (\emph{ennāl}) is native Dravidian, a conditional form of the verb \textquotedblleft{}to say, to sound\textquotedblright{} (DEDR 868, with Tamil \emph{eṉ}): roughly \textquotedblleft{}if one says that\textquotedblright{}.
\end{cousinweb}

\begin{grammarlens}[title={What you've built}]
Five, and a sentence can now go on, hurry, hedge, stop short or turn.
\end{grammarlens}

\subsection*{Your turn}
\begin{itemize}
\raggedright
  \item \emph{Say it:} \emph{ennāl}
  \item \emph{Say it:} it once more, slowly
  \item \emph{Say it:} all five in order — \emph{pinne}, \emph{uṭane}, \emph{cilappōḷ}, \emph{mātraṁ}, \emph{ennāl}
  \item \emph{Recall:} say \emph{muṟaṁ}
\end{itemize}

\begin{tcolorbox}[breakable,colback=teal!4,colframe=teal!35!black,title={Before you move on}]
What does \ml{എന്നാൽ} mean? (\textquotedblleft{}but, on the other hand\textquotedblright{}.) Say it once more.
\end{tcolorbox}
